\documentclass[10pt,a4paper]{amsart}
\usepackage{amsmath,amssymb,amsthm,graphicx}
% Retained preamble for the Stacks Project chapter snapshot in this directory.
% The chapter's own preamble.tex is a copy of the shared Stacks Project preamble; the block
% below is sliced MECHANICALLY, line for line, from it. Nothing was added, removed or
% redefined; the wrapper class, the xr-hyper/hyperref wiring and the externaldocument list
% are omitted because the assembler supplies the document class.
%
% lines taken from preamble.tex: 3, 5, 9, 11, 12, 156, 157, 158, ... (51 lines total)
%
% package coverage restored from the chapter's own preamble (2026-10-13): xy, xy-2cell

\usepackage[all]{xy}
\xyoption{2cell}
\UseAllTwocells

\newenvironment{reference}{\comment}{\endcomment}
\newenvironment{slogan}{\comment}{\endcomment}
\newenvironment{history}{\comment}{\endcomment}
\theoremstyle{plain}
\newtheorem{theorem}[subsection]{Theorem}
\newtheorem{proposition}[subsection]{Proposition}
\newtheorem{lemma}[subsection]{Lemma}
\theoremstyle{definition}
\newtheorem{definition}[subsection]{Definition}
\newtheorem{example}[subsection]{Example}
\newtheorem{exercise}[subsection]{Exercise}
\newtheorem{situation}[subsection]{Situation}
\theoremstyle{remark}
\newtheorem{remark}[subsection]{Remark}
\newtheorem{remarks}[subsection]{Remarks}
\numberwithin{equation}{subsection}
\def\lim{\mathop{\mathrm{lim}}\nolimits}
\def\colim{\mathop{\mathrm{colim}}\nolimits}
\def\Spec{\mathop{\mathrm{Spec}}}
\def\Hom{\mathop{\mathrm{Hom}}\nolimits}
\def\Ext{\mathop{\mathrm{Ext}}\nolimits}
\def\SheafHom{\mathop{\mathcal{H}\!\mathit{om}}\nolimits}
\def\SheafExt{\mathop{\mathcal{E}\!\mathit{xt}}\nolimits}
\def\Sch{\mathit{Sch}}
\def\Mor{\mathop{\mathrm{Mor}}\nolimits}
\def\Ob{\mathop{\mathrm{Ob}}\nolimits}
\def\Sh{\mathop{\mathit{Sh}}\nolimits}
\def\NL{\mathop{N\!L}\nolimits}
\def\CH{\mathop{\mathrm{CH}}\nolimits}
\def\proetale{{pro\text{-}\acute{e}tale}}
\def\etale{{\acute{e}tale}}
\def\QCoh{\mathit{QCoh}}
\def\Ker{\mathop{\mathrm{Ker}}}
\def\Im{\mathop{\mathrm{Im}}}
\def\Coker{\mathop{\mathrm{Coker}}}
\def\Coim{\mathop{\mathrm{Coim}}}
\DeclareMathSymbol{\boxtimes}{\mathbin}{AMSa}{"02}
\def\QCohstack{\mathcal{QC}\!\mathit{oh}}
\def\Cohstack{\mathcal{C}\!\mathit{oh}}
\def\Spacesstack{\mathcal{S}\!\mathit{paces}}
\def\Quotfunctor{\mathrm{Quot}}
\def\Hilbfunctor{\mathrm{Hilb}}
\def\Curvesstack{\mathcal{C}\!\mathit{urves}}
\def\Polarizedstack{\mathcal{P}\!\mathit{olarized}}
\def\Complexesstack{\mathcal{C}\!\mathit{omplexes}}
\def\Pic{\mathop{\mathrm{Pic}}\nolimits}
\def\Picardstack{\mathcal{P}\!\mathit{ic}}
\def\Picardfunctor{\mathrm{Pic}}
\def\Deformationcategory{\mathcal{D}\!\mathit{ef}}
\usepackage{cleveref}
\crefname{originalequation}{original equation}{original equations}
\Crefname{originalequation}{Original equation}{Original equations}
\makeatletter
\newlabel{equation-commutative-square}{{08QM}{source}{}{sourcepointer-1}{}}
\newlabel{equation-commutative-square@cref}{{[originalequation][0][]08QM}{[1][0][]source}{}{}{}}
\makeatother
\crefname{originalequation}{original equation}{original equations}
\Crefname{originalequation}{Original equation}{Original equations}
\makeatletter
\newlabel{equation-double-square}{{08QN}{source}{}{sourcepointer-2}{}}
\newlabel{equation-double-square@cref}{{[originalequation][0][]08QN}{[1][0][]source}{}{}{}}
\makeatother
\crefname{originalremark}{original remark}{original remarks}
\Crefname{originalremark}{Original remark}{Original remarks}
\makeatletter
\newlabel{remark-resolution}{{08QI}{source}{}{sourcepointer-3}{}}
\newlabel{remark-resolution@cref}{{[originalremark][0][]08QI}{[1][0][]source}{}{}{}}
\makeatother
\crefname{originallemma}{original lemma}{original lemmas}
\Crefname{originallemma}{Original lemma}{Original lemmas}
\makeatletter
\newlabel{sites-cohomology-lemma-get-it-now}{{08QA}{source}{}{sourcepointer-4}{}}
\newlabel{sites-cohomology-lemma-get-it-now@cref}{{[originallemma][0][]08QA}{[1][0][]source}{}{}{}}
\makeatother
\title{Flat base change for the lower shriek of the cotangent setting}
\author{The Stacks Project Authors}
\date{stacks-project:cotangent.tex:local-snapshot-20261012; GNU FDL 1.2 or later}
\begin{document}
\maketitle
\setcounter{subsection}{1}
\noindent\textit{Editorial note.} All delivered mathematical text is the author's own, copied from the pinned Stacks Project chapter file cotangent.tex; no statement or proof step was rewritten, reordered or removed.\par\medskip
\noindent\textit{Editorial note.} The statement refers to two labelled equations of the chapter fixing the commutative and double squares; the proof takes the notation from the resolution remark.\par\medskip
\noindent\textit{Editorial note.} The four references of the counted spans are recorded as non-load-bearing navigation pointers; the argument that consumes them stays inside the counted proof span.\par\medskip
\noindent\textit{Editorial note.} Preamble: the retained preamble of this chapter (sources/stacks-cotangent/retained-preamble.tex, lines 1-63), containing the Xy-pic setup, the comment environments, the theorem declarations and the macros including \textbackslash{}NL.\par\medskip
\noindent\textit{Allowed standard prerequisites (editorial).}\begin{itemize}
\item The resolution remark describing the complexes and squares used. Remark remark-resolution, used at proof lines 1183-1186. Supplies the notation and the reductions taken for granted.
\item The criterion computing the lower shriek from a resolution. Cohomology on Sites, Lemma sites-cohomology-lemma-get-it-now, used at proof lines 1185-1188. Reduces the claim to the quasi-isomorphism statement.
\item The comparison between simplicial algebras and their associated complexes. used at proof lines 1188-1194. Identifies the derived tensor product with the complex of the base change.
\item Freely generated resolutions of the algebra B. used at proof lines 1189-1192. Makes s(P\_bullet) a free resolution, so the derived tensor product is computed as claimed.
\end{itemize}\medskip
\ifdefined\hypertarget\hypertarget{sourcepointer-1}{}\fi
\noindent\textit{Source reference.} Original equation 08QM: focused-seeds/AGP/sources/stacks-snapshot-20261011/tags.txt, source lines 11065--11065. Navigation pointer into the Stacks Project; the standard fact it supplies is applied inside the counted proof span and the argument resting on it stays there.\par\medskip
\ifdefined\hypertarget\hypertarget{sourcepointer-2}{}\fi
\noindent\textit{Source reference.} Original equation 08QN: focused-seeds/AGP/sources/stacks-snapshot-20261011/tags.txt, source lines 11066--11066. Navigation pointer into the Stacks Project; the standard fact it supplies is applied inside the counted proof span and the argument resting on it stays there.\par\medskip
\ifdefined\hypertarget\hypertarget{sourcepointer-3}{}\fi
\noindent\textit{Source reference.} Original remark 08QI: focused-seeds/AGP/sources/stacks-snapshot-20261011/tags.txt, source lines 11061--11061. Navigation pointer into the Stacks Project; the standard fact it supplies is applied inside the counted proof span and the argument resting on it stays there.\par\medskip
\ifdefined\hypertarget\hypertarget{sourcepointer-4}{}\fi
\noindent\textit{Source reference.} Original lemma 08QA: focused-seeds/AGP/sources/stacks-snapshot-20261011/tags.txt, source lines 11053--11053. Navigation pointer into the Stacks Project; the standard fact it supplies is applied inside the counted proof span and the argument resting on it stays there.\par\medskip
\begin{lemma}
\label{lemma-flat-base-change}
Assume (\ref{equation-commutative-square}) induces a quasi-isomorphism
$B \otimes_A^\mathbf{L} A' = B'$. Then, with notation as in
(\ref{equation-double-square}) and
$\mathcal{F}' \in \textit{Ab}(\mathcal{C}')$,
we have $L\pi_!(g^{-1}\mathcal{F}') = L\pi'_!(\mathcal{F}')$.
\end{lemma}

\begin{proof}
We use the results of Remark \ref{remark-resolution} without
further mention. We will apply Cohomology on Sites, Lemma
\ref{sites-cohomology-lemma-get-it-now}. Let $P_\bullet \to B$ be a resolution.
If we can show that $u(P_\bullet) = P_\bullet \otimes_A A' \to B'$
is a quasi-isomorphism, then we are done. The complex of $A$-modules
$s(P_\bullet)$ associated to $P_\bullet$ (viewed as a simplicial $A$-module)
is a free $A$-module resolution of $B$. Namely, $P_n$ is a free $A$-module and
$s(P_\bullet) \to B$ is a quasi-isomorphism. Thus $B \otimes_A^\mathbf{L} A'$
is computed by $s(P_\bullet) \otimes_A A' = s(P_\bullet \otimes_A A')$.
Therefore the assumption of the lemma signifies that
$\epsilon' : P_\bullet \otimes_A A' \to B'$ is a quasi-isomorphism.
\end{proof}

\end{document}
